\section*{NeurIPS Paper Checklist}

\begin{enumerate}

\item {\bf Claims}
    \item[] Question: Do the main claims made in the abstract and introduction accurately reflect the paper's contributions and scope?
    \item[] Answer: \answerYes{}
    \item[] Justification: The abstract and introduction state specific, quantified claims (faithfulness improvement from 0.146 to 0.891, 5 retrieval strategies, 4 model families, 8-mode failure taxonomy) that are directly supported by experimental results in Section~4.

\item {\bf Limitations}
    \item[] Question: Does the paper discuss the limitations of the work performed by the authors?
    \item[] Answer: \answerYes{}
    \item[] Justification: Section~4.9 discusses limitations including cross-encoder domain mismatch, evaluation metric coarseness, fine-tuning format damage, and corpus coverage gaps.

\item {\bf Theory assumptions and proofs}
    \item[] Question: For each theoretical result, does the paper provide the full set of assumptions and a complete (and correct) proof?
    \item[] Answer: \answerNA{}
    \item[] Justification: This paper is an empirical applied research project with no theoretical results or proofs.

    \item {\bf Experimental result reproducibility}
    \item[] Question: Does the paper fully disclose all the information needed to reproduce the main experimental results of the paper to the extent that it affects the main claims and/or conclusions of the paper (regardless of whether the code and data are provided or not)?
    \item[] Answer: \answerYes{}
    \item[] Justification: Section~3 provides full details on corpus construction, chunking parameters (800 tokens, 200 overlap), retriever configurations, model identifiers (OpenRouter format), evaluation metrics, judge model, and question set composition. Section~4 reports all hyperparameters for each experiment.

\item {\bf Open access to data and code}
    \item[] Question: Does the paper provide open access to the data and code, with sufficient instructions to faithfully reproduce the main experimental results, as described in supplemental material?
    \item[] Answer: \answerYes{}
    \item[] Justification: Code is provided via GitHub repository with README and requirements.txt. The corpus is constructed from publicly available government websites with scraping scripts included.

\item {\bf Experimental setting/details}
    \item[] Question: Does the paper specify all the training and test details (e.g., data splits, hyperparameters, how they were chosen, type of optimizer) necessary to understand the results?
    \item[] Answer: \answerYes{}
    \item[] Justification: Section~3.6 specifies fine-tuning details (LoRA rank=16, alpha=32, 3 epochs, 586 samples, Unsloth framework, RTX 5080). Section~3.5 specifies evaluation setup (89 questions, Claude Sonnet 4 judge, binary and per-fact metrics).

\item {\bf Experiment statistical significance}
    \item[] Question: Does the paper report error bars suitably and correctly defined or other appropriate information about the statistical significance of the experiments?
    \item[] Answer: \answerYes{}
    \item[] Justification: Section~4.7 reports multi-run averaging with standard deviation and 95\% confidence intervals from 3 parallel runs (e.g., faithfulness 0.891 $\pm$ 0.024, 95\% CI [0.865, 0.918]). The variance source is identified as LLM generation non-determinism.

\item {\bf Experiments compute resources}
    \item[] Question: For each experiment, does the paper provide sufficient information on the computer resources (type of compute workers, memory, time of execution) needed to reproduce the experiments?
    \item[] Answer: \answerYes{}
    \item[] Justification: Fine-tuning was performed on an RTX 5080 (~8 min). Local model inference used llama-server with GGUF F16 format. Cloud models were accessed via OpenRouter API. Evaluation used 3 cloud workers and 3 judge workers in parallel.

\item {\bf Code of ethics}
    \item[] Question: Does the research conducted in the paper conform, in every respect, with the NeurIPS Code of Ethics \url{https://neurips.cc/public/EthicsGuidelines}?
    \item[] Answer: \answerYes{}
    \item[] Justification: The system provides legal information (not advice), uses only publicly available government data, and includes appropriate disclaimers about its limitations.

\item {\bf Broader impacts}
    \item[] Question: Does the paper discuss both potential positive societal impacts and negative societal impacts of the work performed?
    \item[] Answer: \answerYes{}
    \item[] Justification: The introduction discusses the positive impact of improving access to legal information for renters. The paper explicitly distinguishes legal information from legal advice and notes hallucination risks (Section~4.9). The failure taxonomy (Section~4.8) characterizes specific risk modes.

\item {\bf Safeguards}
    \item[] Question: Does the paper describe safeguards that have been put in place for responsible release of data or models that have a high risk for misuse (e.g., pre-trained language models, image generators, or scraped datasets)?
    \item[] Answer: \answerNA{}
    \item[] Justification: The system uses existing LLMs via API and a small LoRA adapter trained on public government data. No high-risk models or datasets are released.

\item {\bf Licenses for existing assets}
    \item[] Question: Are the creators or original owners of assets (e.g., code, data, models), used in the paper, properly credited and are the license and terms of use explicitly mentioned and properly respected?
    \item[] Answer: \answerYes{}
    \item[] Justification: All data sources (mass.gov, boston.gov, MassLegalHelp.org, BostonHousing.org, GBLS) are publicly available government resources. Model providers (OpenAI, Meta, Mistral, Alibaba/Qwen, Anthropic) are credited. ChromaDB and HuggingFace models are cited.

\item {\bf New assets}
    \item[] Question: Are new assets introduced in the paper well documented and is the documentation provided alongside the assets?
    \item[] Answer: \answerYes{}
    \item[] Justification: The curated evaluation dataset (89 questions with key facts and source chunk references) and the failure taxonomy are documented in the paper and code repository.

\item {\bf Crowdsourcing and research with human subjects}
    \item[] Question: For crowdsourcing experiments and research with human subjects, does the paper include the full text of instructions given to participants and screenshots, if applicable, as well as details about compensation (if any)?
    \item[] Answer: \answerNA{}
    \item[] Justification: This project does not involve crowdsourcing or human subjects research.

\item {\bf Institutional review board (IRB) approvals or equivalent for research with human subjects}
    \item[] Question: Does the paper describe potential risks incurred by study participants, whether such risks were disclosed to the subjects, and whether Institutional Review Board (IRB) approvals (or an equivalent approval/review based on the requirements of your country or institution) were obtained?
    \item[] Answer: \answerNA{}
    \item[] Justification: This project does not involve human subjects research.

\item {\bf Declaration of LLM usage}
    \item[] Question: Does the paper describe the usage of LLMs if it is an important, original, or non-standard component of the core methods in this research? Note that if the LLM is used only for writing, editing, or formatting purposes and does \emph{not} impact the core methodology, scientific rigor, or originality of the research, declaration is not required.
    \item[] Answer: \answerYes{}
    \item[] Justification: LLMs are the core component of this research---used as both the response generators (GPT-4o, Llama 3.3 70B, Mistral Small 24B, Qwen3 4B) and the evaluation judge (Claude Sonnet 4). All model identifiers, API providers, and evaluation prompts are fully specified.

\end{enumerate}
